\chapter{結論與未來研究}\label{chap:conclusion}
\begingroup
\setlength{\emergencystretch}{2em}

\section{研究結論}

本研究以具有時間資訊的不平衡分類為對象，在美國企業年度財務資料上比較歷史與近期模型的運用方式。主要分析採用由最後觀測年度回溯重建的事件代理標籤，並以公司最終狀態標籤作為對照。依三項研究問題，結論整理如下。

第一，近期資料的優勢依分類器而異。描述性比較中，XGBoost 與 LightGBM 以近期資料模型的平均 AUC 較高，Random Forest 以合併重訓最高，Logistic Regression 三種策略相近。由於資料期間與筆數同時變動，此結果不單獨證明新近程度的效果，也未對四種分類器作跨設定的顯著性或等效檢定。

第二，在固定的 Old／New 模型池中，New3 與驗證選定靜態集成 V-static 相對 Equal6 的 AUC 增加，皆通過以公司為單位的配對檢定與 165 項 Holm 校正。New3 的平均 AUC 最高；V-static 的平均 F1-score 最高，但其相對 Equal6 的 F1-score 增加未通過校正。這些結果支持將單一期間與全部等權集成保留為控制組，檢查混合歷史與近期模型的效益；本研究未將成員品質與整合規則分開控制，不能比較兩者的相對重要性，也不能由共同基準檢定推論 New3 與 V-static 彼此顯著不同。

第三，特徵選取沒有帶來一致改善。48 項 FS 與同方法 no-FS 的 AUC 比較中，沒有校正後顯著增加；48 項 F1-score 比較亦皆未通過校正，不能將此解釋為等效。SHAP 50\% 的 V-static AUC 下降通過校正，但尾端重抽次數很少，需保留 Monte Carlo 精度限制。Recall 的變化依方法與門檻而異。FS 將特徵由 18 個減至 9 或 14 個，但本研究未證明這能穩定改善少數類辨識或節省實際運算成本。

此外，調參後的單一分類器有接近或高於集成的描述性分數，而集成池成員使用預設參數，兩者驗證資料亦不同，因此尚不足以判斷集成本身是否優於充分調參的單一模型。改用公司最終狀態標籤時，V-static 與 New3 的平均 AUC 順序改變，顯示結果的解讀依賴目標定義；兩種標籤的原始分數與顯著性不直接互相替代。

\section{研究限制}

第一，事件標籤以破產公司的最後觀測年度推定下一年度事件，未逐筆核對破產申請日與財報公開時間，故其效度限於回溯重建的代理目標。驗證資料只有 33 筆事件，DSEL 與 Calibration 分別只有 17 與 16 筆，選模與門檻決策的不確定性較大。

第二，兩批公司配對分析皆為事後探索性，固定原有模型、年份切分與驗證決策，未涵蓋訓練樣本、年份配置或跨公司共同衝擊的不確定性。事件標籤使用 9,999 次重抽，尾端次數較少的顯著性判斷可能受 Monte Carlo 誤差影響；未達顯著不表示方法等效。

第三，集成池成員未經調參，Baseline 與集成之間的分數差同時受參數、分類器與驗證資料影響。第四，所有實驗採固定年份切分，測試期間不更新模型，未驗證持續學習或線上適應，也沒有由年份差異單獨證明概念漂移。第五，實證只使用一份公開資料，同一公司可能跨訓練與測試年度出現，結論不延伸至未見公司或其他市場。

\section{未來研究方向}

後續可先增加事件標籤分析的重抽次數，核對尾端精度敏感的結果，並另行規劃具有明確比較家族的 New3 與 V-static 直接比較；新增分析應標示為後續探索，不回頭依顯著性修改模型。其次，以調參後的成員重建模型池，在相同驗證資料下比較集成與單一分類器，並設計能分離成員組成與整合規則的對照。

第三，採用滑動視窗或逐期更新，檢驗不同年份配置下的穩定性。第四，取得可核對的事件日期或使用其他具有明確事件時間的資料，再區分訓練期間出現過與未出現過的個體，評估推廣範圍。最後，若要主張 FS 降低成本，應以相同程序量測訓練與推論時間，並同時報告辨識效能的變化。
\endgroup
